% ============================================================
%  Section 4 -- Organisation équipe et CI/CD
% ============================================================

\chapter{Organisation de l'équipe et synchronisation Git}

\section{Répartition du travail}

Le projet a été conduit par trois personnes avec une répartition par domaines.
Akram a principalement pris en charge l'agent, la vectorisation, la sécurité
expérimentale et le refactoring final. Widad a porté l'interface d'évaluation,
la knowledge base, le moteur de scoring, les tests A/B et la consolidation des
scores. Ayoub a travaillé sur la recherche KB, le serveur et client MCP,
l'intégration Copilot API, le protocole A2A et la preparation de la soutenance.

\plantumllandscape{gantt_equipe}
{Planification Gantt du projet sous PlantUML.}
{fig:gantt-equipe}

Cette répartition a permis de paralléliser les sujets tout en gardant un point
de convergence hebdomadaire : chaque fonctionnalite devait produire soit une
commande CLI, soit un artefact de benchmark, soit un test automatisable.


\section{Méthode de collaboration et pipeline CI/CD}

L'équipe a organisé son travail autour de branches Git dédiées à chaque lot fonctionnel, avec des merges systématiquement validés par des tests unitaires pour éviter les régressions. Le cycle de développement suivait des étapes courtes : choix du ticket, développement sur branche, tests locaux, demande de merge, corrections, puis synchronisation et régénération des artefacts si besoin.

Le pipeline CI/CD garantissait le respect des standards (style, tests, gates), détectait rapidement les conflits entre travaux parallèles et assurait la traçabilité des résultats sur une base testée et reproductible.

\plantumldiagram{ci_cd_pipeline}{Pipeline Git, CI/CD et production des artefacts.}{fig:ci-cd}

